\section{Proofs}\label{app:proofs}

% ---------------------------------------------------------------------------
\subsection{Proof of Proposition~\ref{prop:consistency}}\label{app:proof_consistency}

We verify the conditions of the extremum-estimator consistency theorem \citep[Theorem~2.1]{newey_mcfadden1994}.

\paragraph{Laplace approximation error.}
Let $\ell(\btheta) = \log\int p(\by,\boldsymbol{\eta}\mid\btheta)\,\mathrm{d}\boldsymbol{\eta}$ denote the exact marginal log-likelihood and $\ell_{\mathrm{LA}}(\btheta)$ the Laplace approximation \eqref{eq:laplace}.
By \citet{tierney_kadane1986}, $|\ell_{\mathrm{LA}}(\btheta) - \ell(\btheta)| = O(1)$ uniformly over $\Theta$ (the error is bounded because the conditional density $f(y_i\mid\eta_i,\pi_i;\,\kappa)$ has bounded derivatives in $\eta_i$ under Assumption~\ref{ass:param}).
Dividing by $n$: $|n^{-1}\ell_{\mathrm{LA}}(\btheta) - n^{-1}\ell(\btheta)| = O(n^{-1}) \to 0$.

\paragraph{Uniform convergence.}
The exact marginal log-likelihood satisfies $n^{-1}\ell(\btheta) \to Q(\btheta) \equiv \E[\log p(y_1\mid\btheta)]$ uniformly over $\Theta$, by the uniform law of large numbers.
The key requirement is that the summands are bounded and equicontinuous in $\btheta$, which holds under the compactness of $\Theta$ (Assumption~\ref{ass:param}) and the boundedness of $\bW$ (Assumption~\ref{ass:W}).
Combined with the vanishing Laplace error, $\sup_{\btheta\in\Theta}|n^{-1}\ell_{\mathrm{LA}}(\btheta) - Q(\btheta)| \plim 0$.

\paragraph{Identification.}
By Assumption~\ref{ass:identification}, $Q(\btheta)$ has a unique maximum at $\btheta_0$.

Finally, compactness of $\Theta$, uniform convergence, and identifiable uniqueness satisfy the conditions of the extremum-estimator theorem, giving $\hat{\btheta}_n \plim \btheta_0$.
\hfill$\square$

% ---------------------------------------------------------------------------
\subsection{Proof of Proposition~\ref{prop:normality}}\label{app:proof_normality}

\paragraph{Score expansion.}
By a mean-value expansion around $\btheta_0$:
\[
  \mathbf{0} = n^{-1/2}\frac{\partial\ell_{\mathrm{LA}}}{\partial\btheta}\bigg|_{\hat{\btheta}}
  = n^{-1/2}\frac{\partial\ell_{\mathrm{LA}}}{\partial\btheta}\bigg|_{\btheta_0}
  + n^{-1}\frac{\partial^2\ell_{\mathrm{LA}}}{\partial\btheta\,\partial\btheta'}\bigg|_{\bar{\btheta}}\;
  \sqrt{n}(\hat{\btheta} - \btheta_0).
\]

\paragraph{CLT for the score.}
The Laplace-approximate score $\partial\ell_{\mathrm{LA}}/\partial\btheta$ can be decomposed as a sum of local contributions $s_i(\btheta)$ that depend on unit $i$ and its spatial neighbours.
Under Assumption~\ref{ass:W}, the spatial weights have uniformly bounded row and column sums, so each $s_i$ depends on at most $O(1)$ other units (the neighbourhood is bounded).
This induces a mixing structure that satisfies the CLT for spatially dependent arrays \citep[see][Theorem~3.1]{lee2004}:
\[
  n^{-1/2}\sum_{i=1}^n s_i(\btheta_0) \dlim N\!\bigl(\mathbf{0},\,\boldsymbol{\Omega}_0\bigr),
\]
where $\boldsymbol{\Omega}_0 = \lim_{n\to\infty} \Var\bigl[n^{-1/2}\sum_i s_i(\btheta_0)\bigr]$.

\paragraph{Hessian convergence.}
By the uniform LLN,
\[
  n^{-1}\frac{\partial^2\ell_{\mathrm{LA}}}{\partial\btheta\,\partial\btheta'}\bigg|_{\bar{\btheta}}
  \plim \mathcal{H}_0 \equiv \E\!\left[\frac{\partial^2 \log p(y_1\mid\btheta)}{\partial\btheta\,\partial\btheta'}\bigg|_{\btheta_0}\right],
\]
which is negative definite by Assumption~\ref{ass:identification}.
Under standard regularity (information matrix equality), $\mathcal{H}_0 = -\mathcal{I}(\btheta_0) = -\boldsymbol{\Omega}_0$.

Combining these results and rearranging the score expansion,
\[
  \sqrt{n}(\hat{\btheta} - \btheta_0) = -\mathcal{H}_0^{-1}\;n^{-1/2}\sum_i s_i(\btheta_0) + o_p(1) \dlim N\!\bigl(\mathbf{0},\,\mathcal{I}(\btheta_0)^{-1}\bigr).
\]
\hfill$\square$

% ---------------------------------------------------------------------------
\subsection{Proof of Proposition~\ref{prop:effects}}\label{app:proof_effects}

Under the restricted ZISAR ($\mathbf{g}=\mathbf{0}$), the marginal expectation of $y_i$ evaluated at the conditional mode is
\[
  \E[y_i] \approx (1 - \pi_i)\,\mu_i, \qquad \mu_i = \exp(\eta_i), \quad \pi_i = \logit^{-1}(\mathbf{z}_i'\bgamma).
\]

The latent predictor at the mode satisfies $\hat{\boldsymbol{\eta}} \approx \mathbf{S}(\rho)\bX\bbeta$ (ignoring the random-effect contribution, which integrates out at the mode).
The partial derivative of $\boldsymbol{\eta}$ with respect to $x_{jr}$ is
\[
  \frac{\partial\eta_i}{\partial x_{jr}} = [\mathbf{S}(\rho)]_{ij}\,\beta_r.
\]
The effect on $\E[y_i]$ is
\[
  \frac{\partial\E[y_i]}{\partial x_{jr}} = (1-\pi_i)\,\mu_i\,[\mathbf{S}(\rho)]_{ij}\,\beta_r + \underbrace{(\text{inflation channel})}_{\text{diagonal only}}.
\]
Stacking: $\mathbf{M}_r = \diag((1-\hat{\bpi})\odot\hat{\bmu})\,\mathbf{S}(\rho)\,\beta_r$.

\textbf{(i)} When $\rho \in [0,1)$ and $\bW \geq 0$ entrywise, the power-series $\mathbf{S}(\rho) = \sum_{s=0}^\infty \rho^s\bW^s$ has all non-negative entries.
Since $(1-\pi_i)\mu_i > 0$, $\mathbf{M}_r$ has the same sign as $\beta_r$.

\textbf{(ii)} When $\rho = 0$, $\mathbf{S}(\rho) = \mathbf{I}_n$ and $\mathbf{M}_r$ is diagonal, so all indirect effects (off-diagonal elements) are zero.
For $\rho \neq 0$, $\mathbf{S}(\rho)$ has positive off-diagonal elements (since $\bW^s$ has positive entries for $s \geq 1$ on connected lattices), making indirect effects non-zero.

\textbf{(iii)} Each element of $\mathbf{S}(\rho) = \sum_{s=0}^\infty \rho^s\bW^s$ is monotonically increasing in $|\rho|$ for $\rho \geq 0$ (since all terms $\rho^s[\bW^s]_{ij}$ are non-negative and increasing in $\rho$).
\hfill$\square$

% =============================================================================
\section{Monte Carlo Tables}\label{app:mc}

Tables~\ref{tab:mc_beta}--\ref{tab:mc_kappa} report the results of the experiment of Section~\ref{sec:montecarlo} for three of the five values of~$\rho$; the summaries for all fifteen scenarios are produced by \texttt{mc\_figure.r}.

\begin{table}[htbp]
\centering
\caption{Monte Carlo bias and RMSE of $\hat{\beta}_1$ ($\beta_1^0 = 0.8$) and $\hat{\beta}_2$ ($\beta_2^0 = -0.5$), $R=1{,}000$ replications per scenario. The ZISAR summaries exclude the fits with $\hat\kappa > 100$ and use 810--963 fits per scenario.}\label{tab:mc_beta}
\footnotesize
\setlength{\tabcolsep}{4pt}
\begin{tabular}{ll rrr rrr rrr}
\toprule
& & \multicolumn{3}{c}{$\rho=0$} & \multicolumn{3}{c}{$\rho=0.3$} & \multicolumn{3}{c}{$\rho=0.6$} \\
\cmidrule(lr){3-5}\cmidrule(lr){6-8}\cmidrule(lr){9-11}
& Model & $n{=}100$ & $n{=}225$ & $n{=}400$ & $n{=}100$ & $n{=}225$ & $n{=}400$ & $n{=}100$ & $n{=}225$ & $n{=}400$ \\
\midrule
\multicolumn{11}{l}{\emph{Panel A: $\hat{\beta}_1$}} \\[3pt]
\multirow{3}{*}{Bias}
  & ZISAR    & $-0.019$ & $-0.005$ & $-0.001$ & $-0.013$ & $0.002$ & $0.000$ & $-0.007$ & $-0.004$ & $0.001$ \\
  & SAR-NB   & $-0.009$ & $-0.004$ & $-0.002$ & $-0.010$ & $-0.000$ & $-0.000$ & $-0.004$ & $-0.003$ & $0.002$ \\
  & ZINB     & $0.001$ & $0.001$ & $0.004$ & $0.008$ & $0.019$ & $0.015$ & $0.056$ & $0.057$ & $0.066$ \\
[3pt]
\multirow{3}{*}{RMSE}
  & ZISAR    & $0.188$ & $0.110$ & $0.083$ & $0.173$ & $0.112$ & $0.080$ & $0.163$ & $0.101$ & $0.079$ \\
  & SAR-NB   & $0.184$ & $0.113$ & $0.086$ & $0.177$ & $0.120$ & $0.084$ & $0.177$ & $0.117$ & $0.090$ \\
  & ZINB     & $0.175$ & $0.107$ & $0.083$ & $0.170$ & $0.113$ & $0.082$ & $0.185$ & $0.128$ & $0.111$ \\
\midrule
\multicolumn{11}{l}{\emph{Panel B: $\hat{\beta}_2$}} \\[3pt]
\multirow{3}{*}{Bias}
  & ZISAR    & $0.008$ & $0.000$ & $0.000$ & $0.003$ & $0.008$ & $0.005$ & $0.003$ & $-0.003$ & $0.001$ \\
  & SAR-NB   & $0.002$ & $-0.003$ & $0.001$ & $0.004$ & $0.010$ & $0.004$ & $0.005$ & $-0.001$ & $0.001$ \\
  & ZINB     & $-0.003$ & $-0.006$ & $-0.002$ & $-0.007$ & $-0.001$ & $-0.006$ & $-0.033$ & $-0.040$ & $-0.041$ \\
[3pt]
\multirow{3}{*}{RMSE}
  & ZISAR    & $0.168$ & $0.108$ & $0.078$ & $0.159$ & $0.105$ & $0.073$ & $0.152$ & $0.091$ & $0.072$ \\
  & SAR-NB   & $0.176$ & $0.110$ & $0.080$ & $0.168$ & $0.111$ & $0.080$ & $0.176$ & $0.107$ & $0.084$ \\
  & ZINB     & $0.165$ & $0.107$ & $0.079$ & $0.161$ & $0.104$ & $0.075$ & $0.176$ & $0.116$ & $0.091$ \\
\bottomrule
\end{tabular}
\end{table}

\begin{table}[htbp]
\centering
\caption{Monte Carlo bias and RMSE of $\hat{\rho}$, $R=1{,}000$ replications per scenario. The ZINB does not estimate~$\rho$.}\label{tab:mc_rho}
\footnotesize
\setlength{\tabcolsep}{4pt}
\begin{tabular}{ll rrr rrr rrr}
\toprule
& & \multicolumn{3}{c}{$\rho=0$} & \multicolumn{3}{c}{$\rho=0.3$} & \multicolumn{3}{c}{$\rho=0.6$} \\
\cmidrule(lr){3-5}\cmidrule(lr){6-8}\cmidrule(lr){9-11}
& Model & $n{=}100$ & $n{=}225$ & $n{=}400$ & $n{=}100$ & $n{=}225$ & $n{=}400$ & $n{=}100$ & $n{=}225$ & $n{=}400$ \\
\midrule
\multirow{2}{*}{Bias}
  & SAR-NB   & $0.063$ & $0.011$ & $0.005$ & $-0.160$ & $-0.174$ & $-0.182$ & $-0.467$ & $-0.477$ & $-0.482$ \\
  & ZISAR    & $-0.022$ & $-0.048$ & $-0.029$ & $-0.024$ & $-0.019$ & $-0.002$ & $-0.021$ & $-0.013$ & $-0.002$ \\
[3pt]
\multirow{2}{*}{RMSE}
  & SAR-NB   & $0.253$ & $0.161$ & $0.115$ & $0.244$ & $0.204$ & $0.196$ & $0.477$ & $0.480$ & $0.484$ \\
  & ZISAR    & $0.433$ & $0.311$ & $0.227$ & $0.338$ & $0.228$ & $0.161$ & $0.229$ & $0.127$ & $0.087$ \\
\bottomrule
\end{tabular}
\end{table}

\begin{table}[htbp]
\centering
\caption{Monte Carlo median bias of $\hat{\kappa}$ ($\kappa^0=2$) and median log-likelihood, $R=1{,}000$ replications per scenario.}\label{tab:mc_kappa}
\footnotesize
\setlength{\tabcolsep}{4pt}
\begin{tabular}{ll rrr rrr rrr}
\toprule
& & \multicolumn{3}{c}{$\rho=0$} & \multicolumn{3}{c}{$\rho=0.3$} & \multicolumn{3}{c}{$\rho=0.6$} \\
\cmidrule(lr){3-5}\cmidrule(lr){6-8}\cmidrule(lr){9-11}
& Model & $n{=}100$ & $n{=}225$ & $n{=}400$ & $n{=}100$ & $n{=}225$ & $n{=}400$ & $n{=}100$ & $n{=}225$ & $n{=}400$ \\
\midrule
\multicolumn{11}{l}{\emph{Panel A: Median bias of $\hat{\kappa}$}} \\[3pt]
  & ZISAR    & $-0.452$ & $-0.511$ & $-0.565$ & $-0.403$ & $-0.523$ & $-0.496$ & $-0.422$ & $-0.353$ & $-0.340$ \\
  & SAR-NB   & $-1.402$ & $-1.428$ & $-1.438$ & $-1.418$ & $-1.448$ & $-1.460$ & $-1.493$ & $-1.511$ & $-1.517$ \\
  & ZINB     & $-0.700$ & $-0.787$ & $-0.860$ & $-0.714$ & $-0.870$ & $-0.896$ & $-0.980$ & $-1.069$ & $-1.103$ \\
\midrule
\multicolumn{11}{l}{\emph{Panel B: Median log-likelihood}} \\[3pt]
  & ZISAR    & $-159$ & $-364$ & $-649$ & $-175$ & $-397$ & $-711$ & $-214$ & $-491$ & $-872$ \\
  & SAR-NB   & $-162$ & $-369$ & $-656$ & $-178$ & $-403$ & $-719$ & $-221$ & $-500$ & $-889$ \\
  & ZINB     & $-160$ & $-365$ & $-650$ & $-176$ & $-399$ & $-713$ & $-220$ & $-499$ & $-887$ \\
\bottomrule
\end{tabular}
\end{table}

\clearpage

\section{Spatial Inflation Predictor: Simulation Checks}\label{app:general}

To assess the Laplace-approximate estimator of the general ZISAR (Section~\ref{sec:model_spec}), we simulate from it on regular lattices with row-standardised queen contiguity, 24 replications for each of $n = 900$, 3600 and 10{,}000, with two i.i.d.\ $N(0,1)$ count covariates, one inflation covariate, $\bbeta = (1,\, 0.5,\, -0.3)'$, $\rho = 0.4$, $\sigma = 0.5$, $\kappa = 2$, $\bgamma = (-1,\, 0.8)'$, $\rho_0 = 0.6$ and $\sigma_0 = 1$.
About 40\% of the simulated counts are zero, four fifths of them structural.
All 72 fits converge, none has a positive Laplace log-likelihood, and the largest gradient component at the optimum is below 0.1.

\begin{table}[htbp]
\centering
\caption{General ZISAR estimated by Laplace approximation on simulated lattices: mean (standard deviation) over 24 replications.}\label{tab:general_mc}
\small
\begin{tabular}{lcccc}
\toprule
 & True & $n = 900$ & $n = 3600$ & $n = 10{,}000$ \\
\midrule
$\beta_1$   & $0.5$  & $0.508$ (0.040) & $0.496$ (0.017) & $0.498$ (0.014) \\
$\beta_2$   & $-0.3$ & $-0.289$ (0.043) & $-0.288$ (0.021) & $-0.304$ (0.012) \\
$\rho$      & $0.4$  & $0.406$ (0.094) & $0.428$ (0.064) & $0.416$ (0.035) \\
$\sigma$    & $0.5$  & $0.443$ (0.111) & $0.441$ (0.041) & $0.444$ (0.027) \\
$\gamma_0$  & $-1$   & $-0.853$ (0.121) & $-0.853$ (0.082) & $-0.867$ (0.038) \\
$\gamma_1$  & $0.8$  & $0.669$ (0.091) & $0.693$ (0.060) & $0.693$ (0.030) \\
$\rho_0$    & $0.6$  & $0.748$ (0.377) & $0.853$ (0.056) & $0.870$ (0.030) \\
$\sigma_0$  & $1$    & $0.344$ (0.089) & $0.295$ (0.065) & $0.288$ (0.040) \\
\bottomrule
\end{tabular}
\end{table}

Table~\ref{tab:general_mc} shows that the count process is recovered as well as under the restricted model: $\bbeta$ and $\rho$ are close to their true values, and $\sigma$ is slightly underestimated, the trade-off with $\kappa$ noted in Section~\ref{sec:montecarlo}.
The inflation process is not.
At every sample size $\sigma_0$ is estimated at about 0.3 instead of 1 and $\rho_0$ at 0.75--0.87 instead of 0.6, and the inflation coefficients are attenuated by 13--17\%.
The bias does not shrink with $n$, because the number of random effects grows with the data and each $g_i$ is informed by a single observation that is zero or not, the setting in which the Laplace approximation of binary-response random effects is known to be most biased \citep{joe2008}.
The approximation replaces a rough field with a smoother one of much smaller variance.
With this estimator $\sigma_0$, $\rho_0$ and, to a lesser extent, $\bgamma$ should not be read at face value, and at the national scale the failure is worse, as the next check shows; the reduced-rank field of Section~\ref{sec:model_spec} is the alternative adopted in the paper.

\paragraph{Reduced-rank field at the national scale.}
The second check uses a design at the scale of the application, built on an outcome with many structural zeros: shoplifting, which cannot be recorded where there are no shops.
The units and the weights matrix are the 32{,}322 units of Section~\ref{sec:solar_data}.
The count process has six covariates---retail and non-retail employment at the workplace (Business Register and Employment Survey 2018, log of one plus), residents (log), population density (log), the income-deprivation rate and the social-rented share---and the inflation process three: an indicator of no recorded retail employment, retail employment and density.
The parameters are the restricted ZISAR estimates for the shoplifting offences recorded by the police in each unit between July 2016 and June 2019 (\texttt{data.police.uk}; 33.9\% of the units record none), which give $\hat\rho = 0$ and attribute 45\% of the zeros to the structural regime.
From these four datasets are simulated with different fields in the inflation predictor: no field; a field of the reduced-rank form \eqref{eq:rr_g} on knots 10~km apart, with $\rho_0 = 0.9$ and $\sigma_0 = 0.5$; a field constant within the areas of the 38 police forces, drawn from $N(0, 0.7^2)$ for each force, the step field that differences of recording between forces would produce; and the unit-level SAR field of the general model, with $\rho_0 = 0.9$ and $\sigma_0 = 0.5$.
Each dataset receives a restricted fit and then the reduced-rank field on the lattice of knots 10~km apart ($K = 1{,}499$ knots, bisquare radius 15~km, about seven knots per unit), by the profile procedure of Appendix~\ref{app:tmb} (\texttt{crime\_application.r}); Table~\ref{tab:general_sim} reports the results.

\begin{table}[htbp]
\centering
\caption{Reduced-rank inflation field (10~km knots) on data simulated at the national scale: likelihood-ratio statistic against the restricted fit, field estimates and correlation between the predicted and the true field. The last row fits the unit-level field of the general model to the data without a field. Bounds of the optimisation box: $\rho_0 \ge 0$ for the reduced-rank field, $\sigma_0 \ge 0.001$.}\label{tab:general_sim}
\footnotesize
\begin{tabular}{lrrrrrrr}
\toprule
Field in the data & Zeros & Structural & SD of field & LR & $\hat\rho_0$ & $\hat\sigma_0$ & Corr. \\
\midrule
None                               & 10{,}962 & 4{,}900 & 0    & 0.0   & 0     & 0.001 & -- \\
Reduced-rank, $(0.9,\,0.5)$        & 10{,}944 & 4{,}827 & 0.85 & 465.7 & 0.82  & 0.76  & 0.88 \\
Constant within forces, SD 0.7     & 10{,}946 & 4{,}920 & 0.61 & 285.9 & 0.90  & 0.40  & 0.81 \\
Unit-level SAR, $(0.9,\,0.5)$      & 11{,}657 & 5{,}834 & 0.99 & 60.9  & 0     & 1.07  & 0.26 \\
\midrule
None, unit-level field fitted      & 10{,}962 & 4{,}900 & 0    & 2{,}089 & $-0.08$ & 2.47 & -- \\
\bottomrule
\end{tabular}
\end{table}

Without a field the profile of the reduced-rank likelihood is below zero at every grid point and the fit returns $\sigma_0$ at its lower bound: the procedure does not invent a field.
With a field it can represent it recovers it: a likelihood ratio of 466, estimates of $(0.82, 0.76)$ for $(0.9, 0.5)$---the field at the units has a standard deviation of 0.85 once the autoregression is applied---and a correlation of 0.88 between the predicted and the true field; the step field of the forces is recovered as well (0.81), by a smooth field that follows the boundaries.
The unit-level field is detected (likelihood ratio 61) but reconstructed poorly (0.26): most of its variance lies below the knot spacing, where the basis cannot follow it.
The last row shows what the unit-level field of the general model does on the data without a field: the profile of its Laplace likelihood rises by 85 points at $\sigma_0 = 1$ before collapsing, and the optimiser runs to $\sigma_0 = 2.5$ and stops with a false-convergence flag at a gradient of order $10^8$, 2{,}089 likelihood-ratio points above the restricted model, where the true statistic is zero.
This is the failure described in Section~\ref{sec:est_laplace}, and it is why the general model is estimated with the reduced-rank field throughout.

\section{TMB Implementation Details}\label{app:tmb}

The \texttt{TMB} template declares $\boldsymbol{\eta}$ as random effects and uses the SAR prior density:
\[
  -\log p(\boldsymbol{\eta}\mid\bbeta,\rho,\sigma) = \frac{n}{2}\log(2\pi\sigma^2) - \log|\det(\mathbf{I}-\rho\bW)| + \frac{1}{2\sigma^2}\|(\mathbf{I}-\rho\bW)\boldsymbol{\eta} - \bX\bbeta\|^2.
\]
The log-determinant is computed via pre-computed eigenvalues $\lambda_1,\ldots,\lambda_n$ of $\bW$ (passed as data).
For possibly complex eigenvalues ($\lambda_j = a_j + b_j\mathrm{i}$):
\[
  \log|\det(\mathbf{I}-\rho\bW)| = \frac{1}{2}\sum_{j=1}^n \log\bigl[(1-\rho a_j)^2 + (\rho b_j)^2\bigr].
\]
For large $n$ (the default in \texttt{zisar\_fit} when $n > 5000$, or with \texttt{logdet = "sparse"}), the template receives instead the sparse matrices $\mathbf{I}_n$, $\bW+\bW'$ and $\bW'\bW$, assembles $\mathbf{Q}(\rho) = \mathbf{I}_n - \rho(\bW+\bW') + \rho^2\bW'\bW$ and evaluates \eqref{eq:logdet_sparse} as $\tfrac12$\texttt{newton::log\_determinant}$(\mathbf{Q})$, an atomic function of \texttt{TMB}'s \texttt{TMBad} framework that factorises $\mathbf{Q}$ by sparse Cholesky and supplies its derivative with respect to $\rho$.
The same mechanism is used by \texttt{TMB} for Gaussian Markov random fields with parameter-dependent precision.
The sparse matrix--vector product $(\mathbf{I}-\rho\bW)\boldsymbol{\eta}$ is computed via Eigen's sparse algebra, preserving the $O(\mathrm{nnz}(\bW))$ complexity.
The reparametrisation $\rho = \tanh(\tilde{\rho})$ ensures $\rho \in (-1,1)$; for numerical stability the optimiser additionally bounds $\tilde\rho\in[-4,4]$, $\kappa\in[0.01,500]$ and $\sigma\in[0.001,10]$ (Section~\ref{sec:est_algorithm}).
Standard errors on the original scale are obtained by the delta method through \texttt{TMB}'s \texttt{ADREPORT} mechanism.

For the general ZISAR (\texttt{zero\_spatial = "field"} in \texttt{zisar\_fit}) the template declares a second random vector $\mathbf{g}$ with the SAR prior above, $\rho_0$ and $\sigma_0$ in place of $\rho$ and $\sigma$ and no mean term, and adds $\mathbf{g}$ to the inflation predictor; $\rho_0 = \tanh(\tilde\rho_0)$ and $\sigma_0$ have the same bounds as $\rho$ and $\sigma$.
For the reduced-rank field \eqref{eq:rr_g} (\texttt{zero\_spatial = "reduced"}) the template receives the basis $\mathbf{B}$, the knot weights $\bW_K$ and their eigenvalues as data, declares the $K$ knot values $\boldsymbol\alpha$ as random effects with the SAR prior on $\bW_K$---its log-determinant from the eigenvalues, $K$ being small---and adds $\mathbf{B}\boldsymbol\alpha$ to the inflation predictor, with $\tilde\rho_0 \ge 0$.
\texttt{zisar\_basis} builds $\mathbf{B}$ and $\bW_K$ from the unit coordinates, the knot spacing and the radius, and \texttt{zisar\_field\_profile} evaluates the Laplace likelihood of either field on a grid of $(\rho_0, \sigma_0)$ with the other parameters at the restricted estimates and then fits the field from the best grid point.
For the SAR filter on the inflation predictor (\texttt{zero\_spatial = "sar"}) the template computes $\boldsymbol{\zeta} = (\mathbf{I}_n - \rho_0\bW)^{-1}\bZ\bgamma$ with \texttt{newton::Newton}, as the minimiser of $\tfrac12\|(\mathbf{I}_n - \rho_0\bW)\boldsymbol{\zeta} - \bZ\bgamma\|^2$: the solver factorises the sparse Hessian $(\mathbf{I}_n - \rho_0\bW)'(\mathbf{I}_n - \rho_0\bW)$ and supplies the derivatives of $\boldsymbol{\zeta}$ with respect to $\rho_0$ and $\bgamma$ by the implicit function theorem, so no random effects are added.
In both specifications the mixture probabilities are computed on the log scale, $\log P(y_i = 0) = \log\{\exp(\log\pi_i) + \exp(\log(1-\pi_i) + \log h(0;\eta_i,\kappa))\}$ with $\log\pi_i = -\log(1 + e^{-\zeta_i})$ and $\zeta_i$ the inflation predictor, which avoids underflow when $\zeta_i$ is large in absolute value.
